\section{The Celestial Hierarchy}\label{sec:ch}

The Church learned the names and the order of the nine choirs from
Scripture, and learned to read them as one sacred order from a small body
of Greek writings that present themselves as the work of Dionysius the
Areopagite, the member of the Athenian council converted by St.~Paul's
sermon on the Areopagus (Acts 17:34). The corpus contains four treatises,
the \emph{Celestial Hierarchy}, the \emph{Ecclesiastical Hierarchy}, the
\emph{Divine Names}, and the \emph{Mystical Theology}, and ten letters;
the treatises also refer to other works of the author, such as a
``Treatise concerning the Divine Hymns'' (\emph{CH} 7.4) and a ``symbolic
theology'' (\emph{CH} 15.6), which are not extant. The \emph{Celestial
Hierarchy} is addressed ``To my Fellow Presbyter Timothy'' and signed
``Dionysius the Presbyter.''

The writings first appear in dated history in the early sixth century.
The earliest reliable citation is by Severus, patriarch of Antioch
(512--518), in a letter to an abbot John. At the conference held at
Constantinople in 533 the Severian speakers quoted Dionysius for the
doctrine of one nature in Christ, and Hypatius of Ephesus, speaking for
the orthodox, rejected the citation as spurious, observing that Cyril of
Alexandria had never used it (J. Stiglmayr, ``Dionysius the
Pseudo-Areopagite,'' \emph{Catholic Encyclopedia} V, 1909). Sergius of
Resaina (d.~536) translated the works into Syriac at an early date. In
1895 Hugo Koch and Stiglmayr, in two independent studies, showed that
the author draws on the Neoplatonist Proclus; the corpus is therefore
dated to about the year 500, and its author is unknown. Benedict~XVI,
devoting a general audience to him, calls him ``a sixth-century
theologian whose name is unknown and who wrote under the pseudonym of
Dionysius the Areopagite,'' places him after Proclus, ``who died in
Athens in 485,'' and prefers to read the pseudonym as an act of
humility: the author ``did not want to glorify his own name'' but
``truly to serve the Gospel, to create an ecclesial theology''
(Benedict~XVI, General Audience, 14 May 2008). The defence of the
Areopagite authorship that the English translator John Parker still
maintained belongs to the history of the corpus and not to its
authority. What the Church received was the teaching, and the teaching
has been tested by its reception.

That reception in the West is early and continuous. Gregory the Great,
preaching on the ninety-nine sheep and the ten drachmas, invokes the
author by name: \latin{Fertur vero Dionysius Areopagita, antiquus
videlicet et venerabilis Pater, dicere} that it is from the lesser companies
that angels are sent to outward ministry (Gregory, \emph{Hom.\ in Evang.}
34.12). In 827 the Byzantine emperor Michael~II sent the writings among
his gifts to Louis the Pious. Hilduin, abbot of Saint-Denis, had them
translated into Latin, and in his life of the saint identified Denis, the
martyr of Paris, with the Areopagite (J.\,P. Kirsch, ``Hilduin, Abbot of
St-Denis,'' \emph{Catholic Encyclopedia} VII, 1910; Stiglmayr,
\emph{loc.\ cit.}). About 858 John Scotus Eriugena made a new Latin
version at the instance of Charles the Bald, which became the main
channel of Dionysius to the Middle Ages. The masters of Saint-Victor,
Hugh foremost among them, built on his doctrine, and Peter Lombard,
Alexander of Hales, Albert, Thomas, and Bonaventure adopted his theses
(Stiglmayr, \emph{loc.\ cit.}); Benedict~XVI says that he was
``virtually rediscovered in the 13th century, especially by St
Bonaventure.'' Aquinas cites the corpus more than ninety times in the two
treatises on the angels alone (Appendix~\ref{app:ch-in-summa}). Lorenzo
Valla (1407--1457) was the first to voice open doubt of the Areopagite
authorship, and the question stayed open until the studies of 1895.

The treatise has fifteen chapters. Chapters 1--2 establish why the
angels are revealed through symbols, even through unlike ones; chapter 3
defines hierarchy; chapters 4--5 treat the name ``angel''; chapters 6--9
give the nine orders in three triads; chapter 10 summarizes; chapters
11--14 answer particular questions of name, mission, and number; and
chapter 15 interprets the bodily images under which Scripture shows the
angels. Benedict~XVI calls the whole a ``symphony of the cosmos that goes
from the Seraphim to the Angels and Archangels.''

Two matters of vocabulary govern every page. First, the Dionysian triads
are, from the highest: \gk{seraphim}, \gk{cheroubim}, \gk{thronoi};
\gk{kyriotetes}, \gk{dynameis}, \gk{exousiai}; \gk{archai},
\gk{archangeloi}, \gk{angeloi}. Latin theology renders the middle triad
\latin{Dominationes, Virtutes, Potestates}, so that Dominations, Virtues,
and Powers correspond to Parker's ``Lordships,'' ``Powers,'' and
``Authorities.'' Quotations below keep Parker's English; the exposition
uses the Latin names. Second, Gregory the Great gives the same nine
names in a different order in the middle of the scale, placing
Principalities above Powers and Virtues below them (Gregory, \emph{Hom.\
in Evang.} 34.7, 10); Aquinas weighs the two orderings and judges that
``little or no difference exists in reality'' between them (\ST{I q.~108
a.~6 ad~4}; see \S\ref{sec:q108} and Appendix~\ref{app:orderings}). This
section follows the Dionysian order throughout. English quotations are
from John Parker's translation, \emph{The Works of Dionysius the
Areopagite}, Part~II (London and Oxford: James Parker and Co., 1899), by
his section numbers; columns are
those of Migne, \emph{Patrologia Graeca} 3. The
orders are catalogued one by one in Appendix~\ref{app:orders}.

\chchapter{That every divine illumination, whilst going forth lovingly to
the objects of its forethought under various forms, remains simplex. Nor
is this all. It also unifies the things illuminated.}
{\emph{CH} 1; PG~3, col.~120B}
{\ST{I q.~108 a.~1 co.}; \ST{I q.~111 a.~1 co.}}
\label{sec:ch1}

The treatise opens with a verse of St.~James, ``Every good gift and every
perfect gift is from above, and cometh down from the Father of Lights''
(Jas 1:17), and draws from it the law that governs all that follows.
Light proceeds from the Father as a gift, and as it comes to us ``restores
us again gradually as an unifying power, and turns us to the oneness of
our conducting Father, and to a deifying simplicity'' (\emph{CH} 1.1).
The procession of light is therefore also a return: what goes forth from
God draws back to God what it reaches.

The second section turns the principle into a prayer. The author invokes
``Jesus, the Paternal Light, the Real, the True, `which lighteth every man
coming into the world','' through whom we have access to the Father, and
asks to be raised to the illuminations ``handed down by our fathers in
the most sacred Oracles,'' there to ``gaze, as we may, upon the
Hierarchies of the Heavenly Minds manifested by them symbolically for our
instruction'' (\emph{CH} 1.2). The study of the angels begins in Christ
and in the Scriptures, not in speculation. The divine ray is multiplied
in going forth and yet ``remains firmly and solitarily centred within
itself in its unmoved sameness,'' and it raises those who aspire to it
and ``makes them one, after the example of its own unifying Oneness.''
The same section gives the sentence that Aquinas makes his own: ``it is
not possible that the supremely Divine Ray should otherwise illuminate
us, except so far as it is enveloped, for the purpose of instruction, in
variegated sacred veils'' (\emph{CH} 1.2). The veils are the images of
Scripture and the rites of the Church.

The third section applies this to the liturgy. The ``Divine Institution
of sacred Rites'' has depicted the immaterial hierarchies ``in material
figures and bodily compositions,'' because the mind cannot rise to them
``without using the material guidance suitable to itself.'' Visible
beauty is a reflection of invisible comeliness; the sweet odours of
incense are ``emblems of the spiritual distribution''; the lights of the
sanctuary are ``a likeness of the gift of the immaterial enlightenment'';
the ranks of the clergy stand for the harmonious order of heaven; and the
reception of the Eucharist stands for ``the partaking of Jesus''
(\emph{CH} 1.3). The Church's hierarchy is constituted ``as
fellow-ministers'' with the heavenly one, and Scripture describes the
angels ``under sensible likeness'' in order to ``lead us through the
sensible to the intelligible, and from inspired symbols to the simple
sublimities of the Heavenly Hierarchies.''

The chapter thus sets the method of the whole work. Angelology is
received from revelation; it is received through images fitted to human
nature; and its end is not information about higher creatures but our
own union with God, in company with them. Aquinas takes from this
chapter the reason why the human hierarchy differs from the angelic:
men receive divine things ``under sensible signs,'' as Dionysius says
(\ST{I q.~108 a.~1 co.}), and it is through the angels' ministry that
revelation reaches men clothed in the likenesses of sensible things
(\ST{I q.~111 a.~1 co.}).

\chchapter{That Divine and Heavenly things are appropriately revealed,
even through dissimilar symbols.}
{\emph{CH} 2; PG~3, col.~136D}
{Not cited by chapter in \ST{I qq.~50--64} or \ST{I qq.~106--114}.}
\label{sec:ch2}

Before speaking of the orders, Dionysius defends the images through
which Scripture speaks of them. He does not want the reader to think,
``like the vulgar,'' that the heavenly minds are ``many-footed and
many-faced creatures,'' or oxen, lions, and eagles, or that there are
``wheels of fire above the heaven, or material thrones upon which the
Godhead may recline'' (\emph{CH} 2.1). Scripture ``artlessly makes use
of poetic representations'' out of regard for our mode of understanding.
There is also a second reason: it is ``most agreeable to the revealing
Oracles to conceal, through mystical and sacred enigmas,'' and to keep
holy truth ``inaccessible to the multitude,'' for ``it is not every one
that is holy'' (\emph{CH} 2.2).

The heart of the chapter is the distinction of two ways of revelation.
``The method of Divine revelation is twofold; one, indeed, as is natural,
proceeding through likenesses that are similar, and of a sacred
character, but the other, through dissimilar forms, fashioning them into
entire unlikeness and incongruity'' (\emph{CH} 2.3). Scripture calls God
Word, Mind, Essence, Light, and Life; these are more reverent names, yet
they ``fall short of the supremely Divine similitude,'' for God is above
every essence and life. At other times Scripture praises him as
invisible, infinite, and incomprehensible, signifying ``not what it is,
but what it is not.'' From this follows a principle of the whole
Dionysian theology: ``If, then, the negations respecting things Divine
are true, but the affirmations are inharmonious, the revelation as
regards things invisible, through dissimilar representations, is more
appropriate to the hiddenness of things unutterable.''

Unlike images are therefore a safeguard. A beautiful image might lead us
to think that the angels are no more than ``men with the appearance of
light'' in shining garments; an incongruous image refuses to let the mind
rest in it and goads it upward (\emph{CH} 2.3). Nor is any image
worthless, since nothing is ``altogether deprived of participation in the
beautiful,'' and ``all things are very beautiful.'' Hence even passions
may be predicated of the angels when rightly purified. What is called
appetite in beasts signifies in the angels ``their manly style, and their
determined persistence''; what is called lust signifies ``a Divine love of
the immaterial, above expression and thought''; incontinence signifies a
love for divine beauty ``which nothing can repulse'' (\emph{CH} 2.4).

The same law holds of God himself. Scripture names him Sun of
Righteousness and Morning Star, but also fire, water, ointment, and
corner-stone, and even Lion, Panther, Leopard, and Bear, and at last a
worm (\emph{CH} 2.5; Ps 21:7). These dissimilar descriptions keep divine
things from the profane and keep the devout from resting ``in the types
as though they were true.'' The chapter ends by asking Timothy to guard
what he receives, for it is not lawful ``to cast to swine'' the
intelligible pearls.

The chapter has no single article of the \emph{Summa} as its heir, but
its doctrine is everywhere assumed in the treatise on the angels: the
bodily images of Scripture are to be taken as signs of spiritual powers,
and no image is to be taken as the angel's nature. The same principle
appears where Aquinas explains the members of assumed bodies as
signifying the powers of the angels, as ``Dionysius teaches'' (\ST{I
q.~51 a.~3 ad~2}; compare \emph{CH} 15).

\chchapter{What is Hierarchy? and what the use of Hierarchy?}
{\emph{CH} 3; PG~3, col.~164D}
{\ST{I q.~108 a.~1 obj.~2 and co.}; \ST{I q.~108 a.~2 obj.~1};
\ST{I q.~108 a.~4 obj.~1}}
\label{sec:ch3}

The word \gk{hierarchia}, ``sacred rule'' or ``sacred principle,''
appears first in these writings, and this chapter gives its classic
definition: ``Hierarchy is, in my judgment, a sacred order and science
and operation, assimilated, as far as attainable, to the likeness of God,
and conducted to the illuminations granted to it from God, according to
capacity, with a view to the Divine imitation'' (\emph{CH} 3.1). Three
terms carry the definition: order (\gk{taxis}), science (\gk{episteme}),
and operation (\gk{energeia}). A hierarchy is not only a ranking; it is
an ordered body that knows and acts, and it exists to receive divine
light and to become like God. God himself, ``as simple, as good, as
source of initiation,'' is ``altogether free from any dissimilarity'' and
gives his light to each according to fitness.

The second section states the purpose: ``The purpose, then, of Hierarchy
is the assimilation and union, as far as attainable, with God, having Him
Leader of all religious science and operation'' (\emph{CH} 3.2). Its
members become ``Divine images, mirrors most luminous and without flaw,''
receiving the primal light and ``spreading this radiance ungrudgingly to
those after it.'' The highest dignity of each member is ``in becoming a
fellow-worker with God'' (1~Cor 3:9). The whole economy is summed up in
a triple law: ``it is an Hierarchical regulation that some are purified
and that others purify; that some are enlightened and others enlighten;
that some are perfected and others perfect'' (\emph{CH} 3.2). The three
acts, purification, illumination, and perfection (\gk{katharsis},
\gk{photismos}, \gk{teleiosis}), recur throughout the corpus. God is
their source, ``cleansing, illuminating, and perfecting,'' and is himself
``above purification, above light, preeminently perfect.''

The third section explains each act from both sides. Those who are
purified must be freed ``from all dissimilar confusion''; those who are
illuminated must be filled with divine light and led to contemplation;
those who are initiated must receive ``that perfecting science of the
sacred things contemplated.'' Correspondingly those who purify give from
their own purity; those who illuminate, ``as being more luminous
intelligences,'' overflow toward those who are worthy; those who perfect
instruct in the science of holy things. Each rank is thus led ``to the
Divine co-operation, by performing, through grace and God-given power,''
what is in God naturally and supernaturally (\emph{CH} 3.3).

This chapter is the source of the concept that organizes \ST{I
qq.~106--108}. Aquinas cites its definition in the objections and
answers of q.~108 a.~1, and he uses the triple law to deny that there is
a hierarchy among the divine Persons: ``for as Dionysius says (Coel.
Hier. iii): `The hierarchical order is so directed that some be cleansed,
enlightened, and perfected; and that others cleanse, enlighten, and
perfect'','' which far be it from us to say of the Persons (\ST{I q.~108
a.~1 co.}). The phrase ``approaches a resemblance to God, as far as may
be'' serves in the objection that the hierarchies belong to grace and not
to nature (\ST{I q.~108 a.~4 obj.~1}).

\chchapter{What is meant by the appellation ``Angels?''}
{\emph{CH} 4; PG~3, col.~177C}
{\ST{I q.~57 a.~1 co.}; \ST{I q.~57 a.~5 obj.~3}; \ST{I q.~106 a.~3
s.c.}; \ST{I q.~111 a.~1 s.c.}; \ST{I q.~113 a.~2 obj.~2}}
\label{sec:ch4}

Dionysius first places the angels within creation. It was ``through
goodness that the superessential Godhead, having fixed all the essences
of things being, brought them into being'' (\emph{CH} 4.1). All things
participate in God's providence: lifeless things by their being, living
things by his life-giving power, rational and intellectual things by his
wisdom. The angels participate more than all others. By moulding
themselves ``to the Divine imitation,'' they ``naturally have more
ungrudging communications with It,'' receiving ``the primal
illuminations without earthly stain,'' and ``having their whole life
intellectual'' (\emph{CH} 4.2).

From this comes the name. They are ``pre-eminently worthy of the
appellation Angelic, on the ground that the supremely Divine illumination
comes to them at first hand, and, through them, there pass to us
manifestations above us'' (\emph{CH} 4.2). The name \gk{angelos},
``messenger,'' belongs to them because they make known what they receive.
Scripture bears witness: the Law was given ``through the ministration of
Angels'' (Gal 3:19; Acts 7:53), and angels led the fathers before and
after the Law, turning them from error and revealing sacred ordinances,
visions, and prophecies.

The third section meets an objection. If some holy men received divine
manifestations directly, how can the angels be mediators? Dionysius
answers that ``no one hath seen, nor ever shall see'' the hidden being of
God as it is in itself (John 1:18); theophanies were visions fitted to
those who saw them, and ``our illustrious fathers were initiated into
these Divine visions, through the mediation of the Heavenly Powers.''
Even the Law, described as coming from God, came ``through Angels, as
though the Divine regulation were laying down this rule, that, through
the first, the second are brought to the Divine Being'' (\emph{CH}
4.3). This is the Dionysian law of mediation, and it holds within each
hierarchy as well: ``there are first, and middle, and last ranks and
powers, and \dots\ the more divine are instructors and conductors of the
less.''

The fourth section applies the law to the mystery of the Incarnation.
The angels ``first were initiated in the Divine mystery of the love of
Jesus towards man,'' and through them the knowledge passed to us.
Gabriel instructed Zachary and revealed to Mary ``how, in her, should be
born the supremely Divine mystery of the unutterable God-formation''; an
angel instructed Joseph; another brought the good news to the shepherds,
and ``a multitude of the Heavenly Host'' sang the doxology (\emph{CH}
4.4; Luke 1--2; Matt 1:20). Even Jesus, ``the superessential Cause of the
super-heavenly Beings,'' when he came among us did not overstep the order
proper to man but received the Father's dispositions through angels, and
was himself called ``Angel of Great Counsel'' (Isa 9:6, Septuagint), since
what he heard from the Father he announced to us.

Aquinas uses the chapter for the angels' nearness to God, from which he
argues that material things pre-exist in them ``more simply and less
materially'' (\ST{I q.~57 a.~1 co.}), and for the principle that
``inferior things are led to God by the superior'' (\ST{I q.~106 a.~3
s.c.}). The revelation of divine things to men ``through the ministry of
the angels'' is the \emph{sed contra} of the question whether angels
enlighten men (\ST{I q.~111 a.~1 s.c.}).

\chchapter{For what reason all the Heavenly Beings are called, in common,
Angels.}
{\emph{CH} 5; PG~3, col.~196B}
{\ST{I q.~108 a.~5 ad~1}; \ST{I q.~113 a.~3 s.c.}}
\label{sec:ch5}

The short fifth chapter, undivided in Parker, asks why ``Angel'' names
both all the heavenly beings together and, in precise speech, only the
last order. Scripture calls them all angels; but ``when they come to a
more accurate description of the supermundane orders, they name
exclusively, `angelic rank,' that which completes the full tale of the
Divine and Heavenly Hosts,'' and ranges above it the Archangels,
Principalities, Powers, Virtues, and the orders higher still (\emph{CH}
5).

The answer rests on a rule that runs through the whole treatise: ``the
superior ranks possess the illuminations and powers of their
subordinates, but the lowest have not the same powers as those who are
above them'' (\emph{CH} 5). The higher orders can therefore be called
angels, since they too ``make known the supremely Divine illumination.''
But the lowest order cannot be called Principalities, Thrones, or
Seraphim, for ``it does not possess the highest powers.'' The name of
the lowest order is common because the lowest property is common; the
names of the higher orders are proper because their excellences are not
shared downward. Dionysius adds a second possible answer, that all the
angelic names are common as regards the communication of all the
heavenly powers toward the divine likeness, and he then turns to examine
``the saintly characteristics set forth respecting each celestial Order
in the Oracles.''

The rule of this chapter is one of the steadiest principles of
Thomist angelology. It explains why the orders are distinct and yet
continuous: each higher order includes what the lower has, and adds
something. Aquinas repeats it in answering why the lowest order bears the
common name. ``The lowest order of angels possess no excellence above the
common manifestation; and therefore it is denominated from manifestation
only; and thus the common name remains as it were proper to the lowest
order, as Dionysius says (Coel. Hier. v)'' (\ST{I q.~108 a.~5 ad~1}). He
states the same rule in general form in answering the question of the
Thrones: ``the excellence of the inferior is contained in the superior,
but not conversely'' (\ST{I q.~108 a.~5 ad~6}).

The chapter also grounds the teaching on guardian angels. Since the
angels properly so called are the order nearest to men, Aquinas can hold
on Dionysius's authority that men are guarded by angels ``who belong to
the lowest order'' (\ST{I q.~113 a.~3 s.c.}, citing \emph{CH} 5 and 9;
see \S\ref{sec:q113}). The name ``angel,'' which in Scripture covers
Gabriel and Michael, the Seraph of Isaiah's vision, and the choirs
before the throne, thus keeps a double sense. It is the common name of
all who carry God's light, and the proper name of those who carry it
last and nearest to us.

\chchapter{Which is the first Order of the Heavenly Beings? which the
middle? and which the last?}
{\emph{CH} 6; PG~3, col.~200C}
{\ST{I q.~56 a.~1 obj.~1}; \ST{I q.~63 a.~7 obj.~1}; \ST{I q.~108 a.~1
s.c.}; \ST{I q.~108 a.~2 co.}; \ST{I q.~108 a.~3 co.}}
\label{sec:ch6}

The sixth chapter, one paragraph in Parker and two sections in Migne,
gives the division of the nine orders into three triads. It opens with
a confession of the limits of our knowledge. How many the orders are,
and how they are ranked, ``the deifying Author of their consecration
alone distinctly knows''; the angels themselves know ``their own proper
powers and illuminations''; and we know only ``so far as the Godhead has
revealed to us, through them'' (\emph{CH} 6). Dionysius therefore
promises to ``utter nothing as from ourselves,'' but to set forth what
the prophets have seen.

Then comes the fundamental sentence: ``The Word of God has designated
the whole Heavenly Beings as nine, by appellations, which shew their
functions. These our Divine Initiator divides into three threefold
Orders'' (\emph{CH} 6). The number nine is not an invention; it is
gathered from the names Scripture gives, and the triple division is the
received interpretation of those names.

The first triad is always around God and ``united closely and
immediately, to Him, before all the rest.'' It comprises ``the most Holy
Thrones, and the many-eyed and many-winged hosts, named in the Hebrew
tongue Cherubim and Seraphim,'' established ``immediately around God,
with a nearness superior to all.'' This triad is ``one, and of equal
rank, and really first Hierarchy, than which there is not another more
Godlike or immediately nearer to the earliest illuminations of the
Godhead.'' The second is ``composed of the Authorities, and Lordships, and
Powers,'' that is, the Powers, Dominations, and Virtues of the Latin
names; the third, ``as respects the lowest of the Heavenly Hierarchies,''
is ``the Order of the Angels and Archangels and Principalities''
(\emph{CH} 6). In this chapter the names are listed from below within
each triad; the proper sequence from the highest is fixed in chapters
7--9.

The chapter sets the architecture of the whole angelic order. Each triad
is a hierarchy; each hierarchy contains three orders; and each order,
as chapter 10 will add, has its own first, middle, and last. The scheme
does not multiply beings for the sake of symmetry; it expresses the law
of chapter 4 that the divine light passes from first to second to last,
and it gives Scripture's scattered names their place in that one
procession.

Aquinas takes the chapter as the \emph{sed contra} for the plurality of
hierarchies: ``Dionysius (Coel. Hier. vi) distinguishes three hierarchies
of angels'' (\ST{I q.~108 a.~1 s.c.}); and ``in every hierarchy Dionysius
places three orders (Coel. Hier. vi)'' (\ST{I q.~108 a.~2 co.}). From
its opening confession he draws the principle that ``our knowledge of the
angels is imperfect, as Dionysius says (Coel. Hier. vi),'' so that we
distinguish their offices only in a general way (\ST{I q.~108 a.~3 co.}).
The same words, that the angels know their own powers, are turned by an
objector into the claim that they do not know their own powers, which
Aquinas answers (\ST{I q.~56 a.~1 obj.~1}). The ordering of the triads is
expounded in \S\ref{sec:q108}.

\chchapter{Concerning the Seraphim and Cherubim and Thrones, and
concerning their first Hierarchy.}
{\emph{CH} 7; PG~3, col.~205B}
{\ST{I q.~57 a.~5 s.c.}; \ST{I q.~106 a.~1 co.}; \ST{I q.~107 a.~2
s.c.}; \ST{I q.~108 a.~5 co., ad~5--6}; \ST{I q.~108 a.~6 s.c., co.};
\ST{I q.~112 a.~4 co.}}
\label{sec:ch7}

Every name of the heavenly minds, Dionysius affirms, ``denotes the
Godlike characteristic of each''; and ``those who know Hebrew affirm,
that the holy designation of the Seraphim denotes either that they are
kindling or burning; and that of Cherubim, a fulness of knowledge or
stream of wisdom'' (\emph{CH} 7.1). The first triad is ministered ``by
the most exalted Beings,'' established immediately around God, to whom
the first divine manifestations pass ``as being nearest.''

The three names are then opened one by one. The name Seraphim teaches
``their ever moving around things Divine, and constancy, and warmth, and
keenness,'' their power to kindle the lower orders ``to the same heat,''
to purify ``through fire and burnt-offering,'' and a light that ``destroys
and dispels every kind of obscure darkness.'' The name Cherubim denotes
``their knowledge and their vision of God, and their readiness to receive
the highest gift of light,'' and their ``ungrudging communication to those
next to them, by the stream of the given wisdom.'' The name Thrones
denotes ``their manifest exaltation above every grovelling inferiority,''
their ``invariable and firmly-fixed settlement around the veritable
Highest,'' and ``their bearing God'' (\emph{CH} 7.1).

The second section describes the hierarchy of the triad. Its members are
pure, not as cleansed from stains, but as ``far exalted above every stain
of remissness''; they are ``never liable to fall, and always unmoved.''
They are contemplative, not through sensible symbols, but ``filled with
all kinds of immaterial knowledge of higher light,'' and ``deemed worthy
of the Communion with Jesus'' in first participation of his
illuminations. They are perfected, not by analytic science, but by ``a
first and pre-eminent deification.'' They receive all this ``not through
other holy Beings, but being ministered from the very Godhead''
(\emph{CH} 7.2).

The third section shows from Scripture that even the highest learn. Some
angels are taught by higher ones that the one ``raised to Heaven as Man,
is Lord of the Heavenly Powers and King of Glory'' (Ps 23:10); others
question Jesus himself, ``Wherefore are Thy garments red?'' (Isa 63:2),
and he teaches them immediately. Dionysius marvels that they first raise
the question ``among themselves,'' not anticipating the illumination.
For the first triad, ``the reception of the supremely Divine Science is,
both purification, and enlightenment, and perfecting'' (\emph{CH} 7.3).
The fourth section describes the triad as the rank that ``dances round
His eternal knowledge in the most exalted ever-moving stability,'' and
recalls its hymns: ``Blessed is the glory of the Lord from His place''
(Ezek 3:12) and ``Holy, Holy, Holy, Lord of Sabaoth, the whole earth is
full of His glory'' (Isa 6:3), treated at length in the lost ``Treatise
concerning the Divine Hymns'' (\emph{CH} 7.4).

No chapter is cited more often in the \emph{Summa}. Aquinas takes from
it the rule that ``the proper name of each order expresses its
property'' (\ST{I q.~108 a.~5 co.}), the exposition of Seraphim from the
properties of fire and of Cherubim as ``fulness of knowledge'' (ad~5),
and of Thrones by their relation to material seats (ad~6). Its sequence,
Seraphim, Cherubim, Thrones, is the \emph{sed contra} of q.~108 a.~6.
The questioning angels of Isaiah and the Psalm ground the teaching that
angels learn the mysteries of grace (\ST{I q.~57 a.~5 s.c.}) and that
the lower speak to the higher (\ST{I q.~107 a.~2 s.c.}).

\chchapter{Concerning Lordships and Powers and Authorities, and
concerning their middle Hierarchy.}
{\emph{CH} 8; PG~3, col.~237B}
{\ST{I q.~106 a.~1 s.c.}; \ST{I q.~108 a.~5 ad~1--3}; \ST{I q.~112
a.~4 obj.~1, s.c.}}
\label{sec:ch8}

The middle triad comprises the Dominations, Virtues, and Powers,
Parker's ``Lordships,'' ``Powers,'' and ``Authorities.'' Dionysius
approaches them ``with supermundane eyes,'' calling their visions ``truly
terrible,'' and repeats that ``each appellation of the Beings above us
manifests their God-imitating characteristics of the Divine Likeness''
(\emph{CH} 8.1).

The Dominations (\gk{kyriotetes}) are named for ``a certain unslavish
elevation, free from all grovelling subserviency,'' superior to every
kind of servility, ``ever aspiring to the true Lordship, and source of
Lordship,'' and moulding themselves and those after them to its lordly
bearing. The Virtues (\gk{dynameis}) are named for ``a certain
courageous and unflinching virility'' in all their godlike energies, not
weak for the reception of any illumination, ``unflinchingly looking to
the superessential and powerful-making power,'' and passing on power to
those next in degree. The Powers (\gk{exousiai}), ``of the same rank as
the Divine Lordships and Powers,'' are named for ``the beautiful and
unconfused good order, with regard to the Divine receptions,'' an
authority ``not using the authoritative powers imperiously for base
purposes,'' but leading those after it ``benignly'' (\emph{CH} 8.1). The
middle triad is purified, illuminated, and perfected ``by the Divine
illuminations vouchsafed to it at second hand, through the first
Hierarchical Order.''

The second section explains why mediated light is dimmer. ``The fulness
of Divine things manifested directly to ourselves is more perfecting
than the Divine contemplations imparted through others''; hence the first
minds are called ``perfecting, and illuminating, and purifying Powers of
the subordinate.'' This is ``divinely fixed absolutely by the Divine
source of order that, through the first, the second partake of the
supremely Divine illuminations'' (\emph{CH} 8.2). Scripture shows it in
many places. In Zechariah one angel, near God, learns the comforting
words and a lower angel comes to meet him to receive them (Zech 1--2).
In Ezekiel the man clothed in linen, ``bound about the loins with a
sapphire,'' receives the judgment on Jerusalem first and the angels with
the battle-axes are instructed from him (Ezek 9). The Cherub casts fire
into the hands of him who wears the sacred vestment (Ezek 10:7), and a
voice commands Gabriel, ``Make this man understand the vision'' (Dan
8:16). By imitating this order ``the discipline of our Hierarchy will
have the Angelic comeliness, as it were, in reflection.''

Aquinas uses the chapter both for the doctrine of illumination and for
the names. ``The angels of the second hierarchy are cleansed,
enlightened and perfected by the angels of the first hierarchy'' is the
\emph{sed contra} of the question whether one angel enlightens another
(\ST{I q.~106 a.~1 s.c.}). The name ``virtues'' signifies ``a certain
virile and immovable strength''; ``Domination'' means ``a certain
liberty, free from servile condition and common subjection''; and
``Power'' signifies a kind of ordination (\ST{I q.~108 a.~5 ad~1--3}).
Because the Dominations are ``above all subjection,'' Aquinas holds with
Dionysius that they are not sent to outward ministry (\ST{I q.~112 a.~4
s.c.}).

\chchapter{Concerning the Principalities, Archangels, and Angels, and
concerning their last Hierarchy.}
{\emph{CH} 9; PG~3, col.~257B}
{\ST{I q.~108 a.~5 ad~3--4}; \ST{I q.~108 a.~6 obj.~4}; \ST{I q.~113
a.~3 s.c.}}
\label{sec:ch9}

The last triad ``completes the Angelic Hierarchies'' and is divided
``into the Godlike Principalities, Archangels, and Angels'' (\emph{CH}
9.1). The name of the Principalities (\gk{archai}) manifests ``their
princely and leading function, after the Divine example,'' their being
``wholly turned to the super-princely Prince, and leading others in
princely fashion.''

The Archangels hold the middle place. Since ``there is not a Hierarchy
which does not possess first and middle and last powers, the holy order
of Archangels occupies the middle position,'' sharing with the
Principalities their turning in princely fashion toward God, and with
the Angels their office of messengers, ``receiving hierarchically the
Divine illuminations from the first powers, and announcing the same to
the Angels'' (\emph{CH} 9.2). The Angels complete the whole series; they
are ``more properly named Angels by us than those of higher degree,
because their Hierarchy is occupied with the more manifest, and is more
particularly concerned with the things of the world.'' The first triad
directs the second hiddenly; the second leads the third ``more clearly
indeed than the first Hierarchy, but more hiddenly than the Order after
it''; and the third ``presides, through each other, over the Hierarchies
amongst men.''

Here the treatise gives its teaching on the angels of the nations.
Scripture ``has assigned our Hierarchy to Angels, by naming Michael as
Ruler of the Jewish people, and others over other nations. For the Most
High established borders of nations according to number of Angels of
God'' (\emph{CH} 9.2; Dan 10:13, 21; Deut 32:8, Septuagint). If only the
Hebrews were led to the true God, the blame lies not with the guiding
angels but with the nations, which ``by their own declension, fell away''
through self-conceit and irrational veneration (\emph{CH} 9.3). The
divine ray is always one and overflowing; the inaptitude of those who
receive it makes their share small or great. Melchizedek, a hierarch
``most dear to God'' outside Israel, shows that the true God had
worshippers among the nations.

The fourth section adds the angels who presided over Egypt and Babylon,
through whom visions were given to Pharaoh and Nebuchadnezzar, with
Joseph and Daniel appointed to interpret them. ``For there is one Prince
and Providence over all.'' We must never think ``that the Godhead is
leader of Jews by lot,'' while angels or other gods rule the nations
independently. The providence of the Most High assigned all mankind
``savingly, to the directing conduct of their own Angels''; Israel,
almost alone, turned to the true Lord and ``became Lord's portion''
(Deut 32:9), and Michael ``became leader of the (Jewish) people''
(\emph{CH} 9.4).

Aquinas takes from this chapter that the name Principalities signifies
``one who leads in a sacred order,'' and that the Archangels stand
``between the `Principalities' and the `Angels','' being ``princes as
regards the `Angels,' and angels as regards the Principalities'' (\ST{I
q.~108 a.~5 ad~3--4}). The same place in the scale is set against
Gregory's order in \ST{I q.~108 a.~6 obj.~4}. The doctrine of guardians
from the lowest order is argued from this chapter with chapter 5 (\ST{I
q.~113 a.~3 s.c.}).

\chchapter{A Repetition and Summary of the Angelic discipline.}
{\emph{CH} 10; PG~3, col.~272D}
{\ST{I q.~50 a.~4 s.c.}; \ST{I q.~108 a.~3 ad~1}; \ST{I q.~113 a.~2
obj.~3}; \ST{I q.~113 a.~3 ad~1}}
\label{sec:ch10}

The tenth chapter gathers the preceding teaching into a single law. The
first order, ``ministered by the perfecting illumination through its
immediate elevation to it,'' is purified, illuminated, and perfected ``by
a gift of light from the Godhead, more hidden and more manifest''; more
hidden because more intelligible, simple, and unifying, more manifest
because it is ``a first gift and a first manifestation'' (\emph{CH}
10.1). From the first order the second is led, from the second the third,
and from the third ``our Hierarchy,'' to ``the super-original Origin and
End of all good order, according to the self-same law of well-ordered
regularity, in Divine harmony and proportion.''

The second section states the principle of interpretation: ``all Angels
are interpreters of those above them, the most reverend, indeed, of God,
Who moves them, and the rest, in due degree, of those who have been moved
by God'' (\emph{CH} 10.2). Every hierarchy is distributed ``into first,
and middle, and last Powers,'' and each division within it is ordered by
the same harmony. Dionysius finds this even within the highest order:
``the theologians say that the most Divine Seraphim cry one to another''
(Isa 6:3), ``indicating distinctly, as I think by this, that the first
impart their knowledge of divine things to the second.'' The
\latin{Sanctus} of Isaiah's vision is not only praise; it is also the
first instruction in heaven.

The third section extends the law to every rational mind. ``Each
heavenly and human mind has within itself its own special first, and
middle, and last ranks, and powers,'' according to which it participates
``in the most spotless purification, the most copious light, the
pre-eminent perfection.'' The chapter ends: ``there is nothing that is
self-perfect, or absolutely without need of perfecting, except the really
Self-perfect and preeminently Perfect'' (\emph{CH} 10.3). The
hierarchical order is therefore not only above us but also within us,
since the soul's own powers are ordered to receive the divine light by
degrees.

The summary shows how tightly the Dionysian vision is woven. The nine
orders, the three triads, and the three acts of purification,
illumination, and perfection are one pattern repeated at every scale:
between the triads, among the orders of a triad, among the members of an
order, and within each mind. God alone stands outside the pattern as its
source, and every creature within it is at once receiver and, according
to its rank, giver.

Aquinas uses the chapter above all for the gradation within a single
order. ``Dionysius says (Coel. Hier. x) that in one and the same order of
angels there are those who are first, middle, and last'' (\ST{I q.~108
a.~3 ad~1}). The same text, cited as \emph{Hier.\ Ang.}\ x, is the
\emph{sed contra} for the thesis that angels of one order differ in
species (\ST{I q.~50 a.~4 s.c.}; see \S\ref{sec:q50}), and it returns in
the question of guardians, where Aquinas holds that ``in each order there
are first, middle, and last'' and that the greater angels probably guard
those chosen by God for greater glory (\ST{I q.~113 a.~3 ad~1}).

\chchapter{For what reason all the Heavenly Beings, in common, are called
Heavenly Powers.}
{\emph{CH} 11; PG~3, col.~284B}
{\ST{I q.~54 a.~3 s.c.}}
\label{sec:ch11}

Chapter 5 asked why all the heavenly beings are called angels; chapter 11
asks the parallel question why they are all called ``Heavenly Powers''
(\gk{ouraniai dynameis}), since the order of Virtues, Parker's ``Powers,''
is not the last. ``For it is not possible to say, as we may of the Angels,
that the Order of the holy Powers is last of all'' (\emph{CH} 11.1). The
answer again uses the rule of inclusion: ``The Orders of the superior
Beings share in the saintly illumination of the last; but the last in no
wise of the first; and on this account all the Divine Minds are called
Heavenly Powers, but never Seraphim and Thrones and Lordships.'' The
orders below the Virtues, Powers, Principalities, Archangels, and Angels,
``are often in common called by us, in conjunction with the other holy
Beings, Heavenly Powers.''

The second section gives a deeper reason and guards against confusion.
The common name does not mix the properties of the orders. ``Inasmuch as
all the Divine Minds, by the supermundane description given of them, are
distributed into three,---into essence, and power, and energy,''
whenever we speak of them all or some of them as heavenly beings or
heavenly powers, we name them in a general way ``from their essence or
power severally'' (\emph{CH} 11.2). Every angel has a being, a power, and
an operation (\gk{ousia}, \gk{dynamis}, \gk{energeia}); to call all
angels ``powers'' is to name them from the second of these. What is
proper to the higher orders must not be applied to the lower ``so as to
overthrow the unconfused order of the Angelic ranks,'' for ``the
superior Orders possess abundantly the sacred characteristics of the
inferior, but the lowest do not possess the superior completeness of the
more reverend.''

The triad of essence, power, and operation passes from this chapter into
the metaphysics of the angel. It supplies the \emph{sed contra} of the
article asking whether the angel's power of intelligence is its essence:
``Dionysius says (Coel. Hier. xi) that `the angels are divided into
substance, power, and operation.' Therefore substance, power, and
operation, are all distinct in them'' (\ST{I q.~54 a.~3 s.c.}). Aquinas
concludes that ``neither in an angel nor in any creature, is the power
or operative faculty the same as its essence'' (\ST{I q.~54 a.~3 co.};
see \S\ref{sec:q54}). A passing remark in Dionysius on a question of names
thus becomes, in the \emph{Summa}, a principle of the real distinction
between the angel's substance and its faculties.

The chapter also explains the language of Scripture and liturgy. The
Psalms bid ``all ye his hosts'' and ``powers'' bless the Lord, and the
Preface joins the Church's praise to that of the heavenly powers. The
common name does not deny the order of the choirs; it names what all of
them share. Each angel, from the Seraph to the guardian, is a power of
God in the sense that it has from God a strength proportioned to its
being and directed to its work.

\chchapter{Why the Hierarchs amongst men are called Angels.}
{\emph{CH} 12; PG~3, col.~292C}
{\ST{I q.~55 a.~3 s.c.}; \ST{I q.~106 a.~4 obj.~1}}
\label{sec:ch12}

The prophet Malachi calls the priest ``the angel of the Lord of hosts''
(Mal 2:7), and the Apocalypse addresses letters to the angels of the
seven churches (Apoc 2--3). Dionysius asks how this can be, since the
lowest orders lack the completeness of the higher: ``for what reason is
our Hierarch named by the Oracles, `Angel of the Sovereign Lord?'''
(\emph{CH} 12.1).

The answer refines the rule of inclusion. The last ``lack the complete
and pre-eminent Power of the more reverend Divisions; for they
participate in the partial and analogous, according to the one
harmonious and binding fellowship of all things'' (\emph{CH} 12.2). The
Cherubim participate in higher wisdom and knowledge; the orders beneath
them participate in wisdom and knowledge too, ``but nevertheless
partially, as compared with them, and in a lower degree.'' Participation
in wisdom is common ``to all the minds which bear the image of God'';
what is not common is the degree of nearness. ``As the first possess
abundantly the saintly characteristics of the inferior, so the last
possess those of the superior, not indeed in the same degree, but
subordinately.'' So the bishop is rightly called angel, ``since he
participates, according to his own capacity, in the messenger
characteristic of the Angels, and elevates himself, as far as attainable
to men, to the likeness of their revealing office.''

The third section goes further. Scripture ``calls gods, both the Heavenly
Beings above us, and the most beloved of God, and holy men amongst us''
(Ps 81:6; John 10:34), although the divine hiddenness is established
above all and no creature can properly be said to be like it, ``except
those intellectual and rational Beings who are entirely and wholly
turned to Its Oneness as far as possible'' and are ``deemed worthy of the
same divine name'' (\emph{CH} 12.3). The name of a higher order may be
given to a lower one by participation; so also the name of God may be
given to creatures who imitate him.

The chapter joins the two hierarchies of the Dionysian corpus. The
bishop who purifies, illuminates, and perfects through the sacraments
does in the Church what the angels do in heaven, and for that reason
receives their name. The bond between them is not a metaphor but a
real participation in one office of making God known.

Aquinas cites the chapter for its teaching on the gradation of angelic
knowledge rather than for the title of bishops. ``Dionysius says (Coel.
Hier. xii) that the higher angels have a more universal knowledge than
the lower'' (\ST{I q.~55 a.~3 s.c.}), which grounds the article on the
universality of the species by which angels understand (see
\S\ref{sec:q55}). An objector uses the same text to argue that a higher
angel does not enlighten a lower concerning all it knows (\ST{I q.~106
a.~4 obj.~1}). Aquinas's own doctrine of participation by excess and by
participation, set out in \ST{I q.~108 a.~5 co.}, gives the terms in
which the Dionysian ``partial and analogous'' sharing is received.

\chchapter{For what reason the Prophet Isaiah is said to have been
purified by the Seraphim.}
{\emph{CH} 13; PG~3, col.~300B}
{\ST{I q.~112 a.~2 s.c., co., ad~2}; \ST{I q.~113 a.~2 obj.~2}}
\label{sec:ch13}

Isaiah says that ``one of the seraphims flew to me, and in his hand was a
live coal'' (Isa 6:6). The difficulty is plain: if the higher orders
are not sent, why is the prophet cleansed not by one of the lower angels
but by one ``enrolled amongst the most reverend Beings''? (\emph{CH}
13.1). Dionysius offers two answers and leaves the reader to weigh them.

The first answer is that the angel sent was one of the angels ``placed
over us as a sacred Minister of the Prophet's cleansing,'' called by the
name of the Seraphim ``on the ground that the removal of the faults
spoken of \dots\ was through fire'' (\emph{CH} 13.2). The name follows
the office: a lower angel who purifies with fire may be called a
Seraph.

The second answer, received from ``another man'' whose defence was ``by
no means foolish,'' is that the angel who formed the vision ``referred
his own cleansing function to God, and after God, to the first working
Hierarchy'' (\emph{CH} 13.3). Divine power reaches all things, yet is
given first to the highest beings and through them to the lower. Two
comparisons illustrate the law. The sun's ray passes easily through the
most transparent matter and grows dim in denser matter; the heat of fire
passes chiefly to things receptive of it and through them to things less
easily heated. So there is for all who are illuminated one source of
light, God, ``by nature, and really, and properly''; but ``by ordinance,
and for Divine imitation, the relatively superior'' is source for each
after it. Hence the lower angels refer every holy energy ``to God indeed
as Cause, but to the first Godlike Minds, as first agents and teachers of
things Divine.''

The fourth section reads the vision itself. Through an angel set over
us the prophet was raised to see the highest beings ``seated \dots\
under God, and with God, and around God,'' and God enthroned above them
all. From the six wings, the feet, the faces, and the ceaseless motion he
learned the Seraphim's power of penetration and their reverence; he was
taught the mysteries of the hymn of praise; and he learned that
purification, like every divine operation, passes to all ``through the
first Beings.'' Therefore the angel ``attributes his own purifying
science and power to God, indeed, as Cause, but to the Seraphim as
first-operating Hierarch,'' just as a bishop is said to purify when he
purifies through his priests (\emph{CH} 13.4). Dionysius ends by asking
Timothy to judge between the answers, or to find a truer one.

Gregory the Great knew this chapter and cited Dionysius for the rule that
the higher choirs do not leave the inner presence of God (Gregory,
\emph{Hom.\ in Evang.} 34.12). Aquinas adopts it through Gregory: ``with
Dionysius (Coel. Hier. xiii) we must say, without any distinction, that
the superior angels are never sent to the external ministry,'' and he
takes the first answer as his own reply to Isaiah: the angel sent ``was
one of the inferior order; but was called a `Seraph,' that is,
`kindling' in an equivocal sense'' (\ST{I q.~112 a.~2 co., ad~2}; see
\S\ref{sec:q112}).

\chchapter{What the traditional number of the Angels signifies.}
{\emph{CH} 14; PG~3, col.~321A}
{\ST{I q.~50 a.~3 co.}; \ST{I q.~112 a.~4 ad~2}}
\label{sec:ch14}

The shortest chapter treats the number of the angels. ``The tradition of
the Oracles concerning the Angels affirms that they are thousand
thousands, and myriad myriads'' (Dan 7:10; Apoc 5:11), ``accumulating and
multiplying, to themselves, the supreme limits of our numbers, and,
through these, shewing clearly, that the ranks of the Heavenly Beings
cannot be numbered by us'' (\emph{CH} 14). The figures of Scripture are
not a census. They take the highest numbers known to us and multiply
them by themselves, to say that the number of the angels exceeds every
count we can make.

Dionysius gives the reason: ``many are the blessed hosts of the
supermundane minds, surpassing the weak and contracted measurement of our
material number.'' They are known precisely only ``by their own
supermundane and heavenly intelligence and science alone,'' which is
given them ``by the supremely Divine and Omniscient Framer of Wisdom,''
the cause and bond and term ``of all created things together'' (\emph{CH}
14). The number of the angels is not unknown to God, who made them and
holds them in being; and it is known to the angels themselves by the
knowledge God gives them. It is unknown to us because our numbering is
material.

The chapter supports a firm tradition of the Church. Daniel sees
``thousands of thousands'' ministering and ``ten thousand times a hundred
thousand'' standing before the Ancient of Days (Dan 7:10), and the Lord
speaks of ``more than twelve legions of angels'' at his command (Matt
26:53). The Fathers read these as signs of an immense multitude, and the
nine orders are therefore nine orders of countless spirits, not nine
spirits.

Aquinas makes the chapter the conclusion of his argument on the
multitude of the angels. Since the perfection of the universe is chiefly
intended by God, the more perfect things are, the more they are created
in excess; the angels, being immaterial, exceed all material multitude.
``This is what Dionysius says (Coel. Hier. xiv): `There are many blessed
armies of the heavenly intelligences, surpassing the weak and limited
reckoning of our material numbers''' (\ST{I q.~50 a.~3 co.}; see
\S\ref{sec:q50}). Corderius's note in Migne on this chapter points to the
same article. In the question of the angels' missions Aquinas returns to
it: Dionysius ``declares that the multitude of angels surpasses all the
multitude of material things,'' and reads the thousands and ten
thousands of Daniel so that those who assist outnumber those who
minister, ``because whatever is better is more intended and more
multiplied by God'' (\ST{I q.~112 a.~4 ad~2}; see \S\ref{sec:q112}). The multitude of angels is
thus common teaching, grounded in Scripture and received in the Church
from Dionysius through the schools.

\chchapter{What are the morphic likenesses of the Angelic Powers? what
the fiery? what the anthromorphic? \dots}
{\emph{CH} 15; PG~3, col.~325D}
{\ST{I q.~51 a.~3 ad~2}; \ST{I q.~106 a.~1 co.}; \ST{I q.~106 a.~4
s.c.}; \ST{I q.~107 a.~5 obj.~3}}
\label{sec:ch15}

The last and longest chapter returns to the images with which the
treatise began, now reading them one by one. Its title lists them: fire,
the human form, eyes, nostrils, ears, mouths, hands, heart, feet, wings,
garments, girdles, rods, spears, winds, clouds, brass, stones, lion, ox,
eagle, horses, rivers, chariots, wheels, and ``the so-called joy of the
Angels.'' The first section sets a key: the same ranks are shown
``sometimes ruling, and, at other times, as being ruled,'' since each
order is ruled by those above and rules those below, and all participate
``in the providential faculty to provide for those below them''
(\emph{CH} 15.1).

Fire comes first because Scripture prefers it: wheels of fire, living
creatures of fire, rivers of flame, fiery thrones, and the Seraphim whose
name means burning. ``The similitude of fire denotes the likeness of the
Heavenly Minds to God in the highest degree'' (\emph{CH} 15.2). Sensible
fire is in everything and unmingled with anything, luminous yet hidden,
``ever moving, self-moving, moving other things,'' present to all
invisibly and undiminished in giving itself; it is therefore apt to
signify the divine energy and the angels who imitate it.

The human form signifies the intellectual faculty, the power of looking
upward, and the ruling of creatures ``by their superior power of mind''
(\emph{CH} 15.3). Each member is then read: sight as ``the most
transparent elevation towards the Divine lights,'' hearing as the
reception of divine inspiration, touch as discrimination, the eyelids as
guarding ``the conceptions which see God,'' the hands as ``the power of
making, and operating, and accomplishing,'' the heart as ``a symbol of
the Godlike life,'' and the feet under the wings as ``the moving and
quickness'' of their progress to divine things. Of the teeth Dionysius
writes that ``each intellectual Being divides and multiplies, by a
provident faculty, the unified conception given to it by the more Divine
for the proportionate elevation of the inferior.''

Garments and instruments follow (\emph{CH} 15.4--5): shining raiment
signifies the divine likeness after the image of fire; the priestly
robe, their leading to mystical visions; girdles, the guarding of their
powers; rods, ``the kingly and directing faculty''; spears and axes, the
discerning of unlike things; measuring instruments, their building and
guiding care. They are called winds for their swift action and clouds
because they are ``filled super-mundanely with the hidden Light''
(15.6). Brass, electrum, and coloured stones signify brilliance and fire
(15.7); the lion signifies the leading and indomitable, the ox the strong
and mature, the eagle the kingly and soaring gaze toward the divine sun,
and horses obedience (15.7--8). The rivers of fire are divine streams,
the chariots the communion of equals, and the winged wheels their
straight advance; the name given to the wheels, ``Gel, Gel,'' ``shews,
according to the Hebrew tongue, revolutions and revelations'' (Ezek
10:13; \emph{CH} 15.9). The angels' joy is no passion but a ``Godlike and
ungrudging rejoicing'' over those who turn to God (Luke 15:10). What
cannot be explained, the treatise concludes, ``we have honoured by
silence.''

Aquinas takes from the chapter his principle of the communication of
knowledge downward: ``Every intellectual substance with provident power
divides and multiplies the uniform knowledge bestowed on it'' (\ST{I
q.~106 a.~1 co.}; compare a.~4 s.c.). The bodily members of assumed
bodies signify angelic powers, ``as Dionysius teaches'' (\ST{I q.~51 a.~3
ad~2}; see \S\ref{sec:q51}).
